Le chlore existe dans la nature sous forme de trois isotopes~: le chlore 35
$(\ce{_{17}^{35}Cl})$, le chlore 36 $(\ce{_{17}^{36}Cl})$ et le chlore 37
$(\ce{_{17}^{37}Cl})$.

Dans les eaux de surface (mers, lacs, \ldots), le chlore 36 est constamment
renouvelé et sa teneur peut être supposée constante.

Dans le cas des eaux de nappes phréatiques, le renouvellement n'existe plus et
la proportion en chlore 36 diminue au cours du temps.

\begin{donnees}
    \begin{table}[H]
        \centering
        \setlength{\arrayrulewidth}{0.8pt}
        \renewcommand{\arraystretch}{1.5}
        \begin{tabularx}{\textwidth}{|p{4cm}|C|C|C|}
            \hline
            Noyau ou particule & Électron & Chlore $\ce{_{17}^{36}Cl}$ &
            Argon $\ce{_{18}^{36}Ar}$ \\
            \hline
            Masse en (\unit{\u}) & $0,000549$ & $35,968312$ & $35,967545$ \\
            \hline
            \multicolumn{3}{|l|}{Constante radioactive du chlore 36~:
            $\lambda=\qty{2.30e-6}{\per\mathrm{ans}}$} &
            \multicolumn{1}{c|}{$\qty{1}{\u}=\qty{931.5}{\U}$} \\
            \hline
        \end{tabularx}
        \begin{tabularx}{\textwidth}{|p{10cm}|C|C|C|}
            \hline
            Noyau & $\ce{_{17}^{35}Cl}$ & $\ce{_{17}^{36}Cl}$ & 
            $\ce{_{17}^{37}Cl}$ \\
            \hline
            \raisebox{0mm}[7mm][4mm]{
                Énergie de liaison par nucléon~:
                $\frac{E_{\ell}}{A}$ ($\unit{\MeV\per\mathrm{nucléon}}$)} &
            $8,5178$ & $8,5196$ & $8,5680$ \\
            \hline
        \end{tabularx}
    \end{table}
\end{donnees}

\begin{enumerate}
    \item Recopier sur votre copie le numéro de la question et écrire la lettre
          correspondante à la proposition vraie.

          La composition du noyau de chlore $\ce{_{17}^{35}Cl}$ est~:

          \begin{minipage}{\linewidth}
              \begin{table}[H]
                  \centering
                  \setlength{\arrayrulewidth}{0.8pt}
                  \renewcommand{\arraystretch}{1.5}
                  \begin{tabularx}{\textwidth}{|p{1cm}|C|}
                      \hline
                      A & 17 protons et 35 neutrons \\
                      \hline
                      B & 18 protons et 17 neutrons \\
                      \hline
                      C & 17 protons et 18 neutrons \\
                      \hline
                      D & 18 protons et 35 neutrons \\
                      \hline
                  \end{tabularx}
              \end{table}
          \end{minipage}
    
    \item Déterminer, en justifiant votre réponse, le noyau le plus stable parmi
          $\ce{_{17}^{35}Cl}$, $\ce{_{17}^{36}Cl}$ et $\ce{_{17}^{37}Cl}$.

    \item Le chlore 36 radioactif, donne en se désintégrant un noyau d'argon
          $\ce{_{18}^{36}Ar}$.
          \begin{enumerate}
              \item Écrire l'équation de désintégration d'un noyau de chlore 36
                    et identifier le type de cette désintégration.
              \item Calculer, en unité (\unit{\MeV}), l'énergie libérée
                    $E_{\text{libérée}}=|\Delta E|$ au cours de la
                    désintégration d'un noyau de chlore 36.
          \end{enumerate}

    \item Un échantillon de volume $V$, des eaux de surface, contient $N_{0}$
          noyaux de chlore 36. Un échantillon de même volume $V$, d'eau issue
          d'une nappe phréatique ne contient que $38 \%$ du nombre de noyaux de
          chlore 36 trouvé dans les eaux de surface.
    
          Déterminer, en unité (ans), l'âge de la nappe phréatique.
\end{enumerate}

